% Generated draft: requires mathematical and source review.
\documentclass[11pt]{article}
\usepackage[margin=25mm]{geometry}
\usepackage{fontspec}
\setmainfont{Latin Modern Roman}
\usepackage{amsmath,amssymb,mathtools,unicode-math}
\setmathfont{Latin Modern Math}
\usepackage{enumitem,needspace,xurl,hyperref}
\setlength{\parindent}{0pt}
\setlength{\parskip}{5pt}
\setlength{\emergencystretch}{3em}
\newcommand{\N}{\mathbb{N}}
\newcommand{\Z}{\mathbb{Z}}
\newcommand{\Q}{\mathbb{Q}}
\newcommand{\R}{\mathbb{R}}
\newcommand{\C}{\mathbb{C}}
\newcommand{\F}{\mathbb{F}}
\newcommand{\A}{\mathbb{A}}
\newcommand{\PP}{\mathbb{P}}
\newcommand{\E}{\mathbb{E}}
\newcommand{\eps}{\varepsilon}
\newcommand{\dd}{\,\mathrm{d}}
\begin{document}
\section*{UnsolvedMath ID: 800596}
\textbf{Source problem code:} catalogue entry
\subsection{Definitions and mathematical statement}
The supplied question concerns the five properties
\[
\mathsf{S}_1(\mathrm{T},\mathrm{T}),\quad
\mathsf{S}_{\mathrm{fin}}(\mathrm{T},\mathrm{T}),\quad
\mathsf{S}_1(\Gamma,\mathrm{T}),\quad
\mathsf{S}_{\mathrm{fin}}(\Gamma,\mathrm{T}),\quad
\mathsf{U}_{\mathrm{fin}}(\mathcal{O},\mathrm{T}),
\]
and their Borel versions. The excerpt does not specify the class of spaces, cover conventions, or meaning of additivity. The following gives a conditional formulation using conventional countable-cover definitions; these conventions require confirmation against the source.

Let \(X\) be a topological space and let \(\N=\{1,2,\ldots\}\). A countable cover of \(X\) is a finite or countably infinite family \(\mathcal U\) of subsets of \(X\) such that \(\bigcup\mathcal U=X\) and \(X\notin\mathcal U\). Let \(\mathcal O(X)\) denote the family of countable open covers under this convention. A cover is large if each point belongs to infinitely many of its members. Define \(\Gamma(X)\) to consist of the infinite members \(\mathcal U\) of \(\mathcal O(X)\) such that
\[
\forall x\in X,\qquad
\bigl|\{U\in\mathcal U:x\notin U\}\bigr|<\infty.
\]
Define \(\mathrm T(X)\) to consist of the large members \(\mathcal U\) of \(\mathcal O(X)\) such that, for every \(x,y\in X\),
\[
\bigl|\{U\in\mathcal U:x\in U,\ y\notin U\}\bigr|<\infty
\quad\text{or}\quad
\bigl|\{U\in\mathcal U:y\in U,\ x\notin U\}\bigr|<\infty.
\]
These are the open \(\gamma\)-covers and open \(\tau\)-covers, respectively.

For families \(\mathscr A(X)\) and \(\mathscr B(X)\) of covers, \(\mathsf S_1(\mathscr A,\mathscr B)\) means that for every sequence \((\mathcal U_n)_{n\in\N}\) from \(\mathscr A(X)\), there exist \(U_n\in\mathcal U_n\) such that
\[
\{U_n:n\in\N\}\in\mathscr B(X).
\]
The property \(\mathsf S_{\mathrm{fin}}(\mathscr A,\mathscr B)\) means that for every such sequence there exist finite subfamilies \(\mathcal F_n\subseteq\mathcal U_n\) such that
\[
\bigcup_{n\in\N}\mathcal F_n\in\mathscr B(X).
\]
The property \(\mathsf U_{\mathrm{fin}}(\mathscr A,\mathscr B)\) means that for every such sequence in which no \(\mathcal U_n\) has a finite subcover, there exist finite subfamilies \(\mathcal F_n\subseteq\mathcal U_n\) such that
\[
\left\{\bigcup\mathcal F_n:n\in\N\right\}\in\mathscr B(X).
\]
Here selected families are treated as sets, so repeated selected subsets do not count as distinct cover members.

For the corresponding Borel formulation, let \(\mathcal B(X)\) be the countable covers by Borel subsets of \(X\), again excluding \(X\) as a member. Let \(\mathcal B_\Gamma(X)\) and \(\mathcal B_{\mathrm T}(X)\) satisfy the same respective pointwise conditions as \(\Gamma(X)\) and \(\mathrm T(X)\). Under this interpretation, the five Borel versions are
\[
\begin{gathered}
\mathsf S_1(\mathcal B_{\mathrm T},\mathcal B_{\mathrm T}),\qquad
\mathsf S_{\mathrm{fin}}(\mathcal B_{\mathrm T},\mathcal B_{\mathrm T}),\\
\mathsf S_1(\mathcal B_\Gamma,\mathcal B_{\mathrm T}),\qquad
\mathsf S_{\mathrm{fin}}(\mathcal B_\Gamma,\mathcal B_{\mathrm T}),\qquad
\mathsf U_{\mathrm{fin}}(\mathcal B,\mathcal B_{\mathrm T}).
\end{gathered}
\]

\textbf{Question, with the unresolved scope made explicit.} Are any of the five listed open-cover properties, or any of their Borel versions, additive? If additivity means preservation under binary unions of subspaces, fix an ambient topological space \(Z\) and an intended class \(\mathscr D\) of its subspaces that is closed under finite unions. For each listed property \(P\), the assertion to be decided is
\[
\forall X,Y\in\mathscr D,\qquad
\bigl(P(X)\land P(Y)\bigr)\Longrightarrow P(X\cup Y),
\]
where each space carries its subspace topology. The source excerpt does not determine \(Z\), \(\mathscr D\), or whether this binary-union interpretation is intended. In particular, preservation under countable unions is a separate possible target and is not asserted to be equivalent. Thus the displayed framework is conditional, and a uniquely specified formulation requires the missing source conventions.
\subsection{Short English statement}
The supplied Problem 24 asks which, if any, of five selection principles involving \(\tau\)-covers, or their Borel versions, are additive. Under a binary-union interpretation, this asks whether the union of two spaces with a given property must also have that property. The intended class of spaces and the precise cover and additivity conventions remain unspecified in the supplied excerpt.
\subsection{Sources}
\begin{itemize}
\item Supplied local UnsolvedMath registry record: permanent ID \(800596\), exact source problem code \texttt{catalogue entry}; upstream ID \(800596\), revision \texttt{d371ee917ac9647fbb06cd6fc529deec6f3dda82}. The registry attributes the problem to Tsaban, with proposed year \(2006\), and labels its status as open. This is a supplied status label, not an independently verified assessment.
\item Tsaban, \emph{Selection Principles and Special Sets of Reals: Open Problems} (2006), Problem 24. Supplied source URL: \url{https://arxiv.org/abs/math/0609601}. The record reports access on 2026-10-16 and extraction by \texttt{arxiv\_latex}; these are supplied provenance details, not a claim of access during this preparation.
\item The supplied title and background associate the problem with the context label \texttt{[AddQuad]}. No bibliographic expansion or explanation of that label is supplied.
\end{itemize}
\end{document}
